\chapter{La expropiación}

\label{cap:expropiacion}

Este capítulo empieza donde empieza todo lo demás.

El 8 de junio de 1972, el gobernador de la provincia de Salta sancionó y promulgó, en un solo acto y sin intervención de ninguna cámara legislativa, una ley que declaraba de utilidad pública seiscientas ochenta y seis hectáreas del departamento La Caldera. La tierra se destinaba a construir un embalse que abastecería de agua potable al sector norte de la ciudad de Salta.

Cincuenta y cuatro años después, ese embalse abastece a la capital. El municipio que lo aloja no tiene red de gas natural ---en el departamento, sólo Vaqueros la tiene---, tiene cobertura cloacal parcial y tiene desigualdad documentada en el acceso al agua potable. En la zona alta de Vaqueros, el otro municipio del departamento, había vecinos que en 2021 llevaban más de quince años esperando el camión municipal dos o tres veces por semana.

Entre esas dos frases está el libro entero.

\section{La ley}

El nombre del embalse proviene de la finca Campo Alegre. Quién era su titular y bajo qué régimen la Provincia tomó esa tierra se responde en la Ley provincial Nº 4486.

\begin{ficha}{Ley provincial Nº 4486}
Sancionada y promulgada el 8 de junio de 1972. Publicada en el Boletín Oficial Nº 9.055 el 19 de junio de 1972. Declara de utilidad pública y sujetas a expropiación cinco fracciones del departamento La Caldera, afectadas por la construcción del \textbf{Embalse de Campo Alegre} y el \textbf{desvío de la Ruta Nacional Nº 9}:
\begin{center}\small
\begin{tabular}{@{}p{4.6cm}p{1.3cm}p{3.4cm}r@{}}
\toprule
Propiedad & Catastro & Titular & Superficie (ha) \\
\midrule
\textbf{Campo Alegre} & 136 & Isaac Fernández & 167,4139 \\
\textbf{Los Porongos} & 1093 & Jorge Washington Álvarez & 1,1720 \\
Fracciones de Loma Bola, Campo de Arrieta y Cañada Ancha & 181 & Francisco Javier Mercado & 1,8036 \\
El Angosto o Angostura & 60 & Jaime Durán & 152,7258 \\
Fracciones de Loma Bola, Campos de Arrieta y Cañada Ancha & 1575 a 1578 & María Celia Mercado de Fernández, Francisco Javier Mercado, Margarita Mercado Vda.\ de Cáceres y Juana Mercado de Sánchez & 363,0718 \\
\midrule
\multicolumn{3}{@{}l}{\textbf{Total expropiado}} & \textbf{686,1871} \\
\bottomrule
\end{tabular}
\end{center}
Planos de expropiación Nros.\ 00136, 00137, 00143, 00144 y 00148 del departamento La Caldera. Firmantes: Spangenberg --- Sosa.
\end{ficha}

\subsection*{Quiénes perdieron la tierra}

Seiscientas ochenta y seis hectáreas, para un embalse cuyo espejo de agua ocupa trescientas veintidós. La diferencia corresponde al perilago, las obras complementarias y el desvío de la ruta.

\textbf{La familia Mercado aportó 364,88 hectáreas: el 53,2 por ciento de todo lo expropiado.} Aparece en dos de los cinco incisos, primero Francisco Javier Mercado a título individual y luego cuatro titulares --- María Celia Mercado de Fernández, Francisco Javier Mercado, Margarita Mercado Vda.\ de Cáceres y Juana Mercado de Sánchez --- sobre catastros correlativos, 1575 a 1578, con títulos registrados a folios correlativos del mismo libro. Es la partición de un dominio familiar entre cuatro herederos, expropiada en bloque.

\begin{cautela}
\textbf{Y ese dominio familiar ya estaba en el archivo, veintiocho años antes.} El Revalúo General de 1944 que transcribe el capítulo \ref{cap:fincas} registra, bajo la partida 80, a «\textbf{Sánchez, Juana Mercado de, y otros}» con una valuación de \$26.400: \textbf{la tercera más alta del departamento}, el 7,5 por ciento de todo el valor publicado, detrás de Urquiza\index{Urquiza, Francisco S.} y de Rauchi.

Es el mismo nombre completo que la Ley 4486 consigna en 1972 entre los cuatro titulares de los catastros 1575 a 1578. \textbf{Como todas las de este libro, es una coincidencia de nombre y no una identificación registral} ---se acredita con el folio, no con el padrón de valuaciones---, y la sucesión pudo conservar el nombre de la causante durante décadas. Lo que la coincidencia aporta, si se confirma, es el eslabón que el capítulo \ref{cap:fincas} pedía: \textbf{la tierra que el Estado expropió para el embalse era, veintiocho años antes, uno de los tres mayores patrimonios registrados del departamento}. Y algo parecido cabe decir de «María Celia Mercado \textbf{de Fernández}» junto a «\textbf{Isaac Fernández}», titular de Campo Alegre en el mismo acto. Las dos preguntas se contestan con las mismas partidas y folios que el apéndice \ref{cap:pedidos} reclama.
\end{cautela}

La finca Campo Alegre propiamente dicha, de Isaac Fernández, aportó 167,41 hectáreas, el 24,4 por ciento. El Angosto o Angostura, de Jaime Durán, 152,73, el 22,3.

\textbf{Primera.} \textit{Los Porongos}, catastro 1093, titular Jorge Washington Álvarez, es una de las fincas que el artículo 6º del proyecto de ordenanza del Plan de Desarrollo Urbano Ambiental enumera cuarenta y tres años después como localidad o paraje integrante del municipio. El capítulo \ref{cap:fincas} planteó esa lista como el mapa de tenencia histórica del departamento y dejó pendiente su reconstrucción dominial. Ésta es la primera pieza: en 1972, Los Porongos tenía titular con nombre, catastro y folio.

\textbf{Segunda.} La estructura que la ley revela --- un dominio familiar dividido entre cuatro titulares del mismo apellido sobre catastros correlativos --- es exactamente la misma que el capítulo \ref{cap:fincas} documentó para la sucesión Serrey\index{Serrey, Carlos} en Vaqueros veinte años después. La diferencia es el destino: en Vaqueros la partición abrió el mercado de suelo; en Campo Alegre, la expropiación.

\subsection*{Cómo se pagó}

El artículo 3º de la ley fija el mecanismo, y es la parte que importa.

La Dirección General de Inmuebles determina la \textbf{evaluación fiscal} de cada fracción. Cumplido eso, la Fiscalía de Gobierno solicita la \textbf{inmediata posesión} depositando esa evaluación fiscal \textbf{incrementada en un 30 por ciento}, conforme el artículo 17 de la Ley 4.172/68. Recién después se realiza el trámite del artículo 13 de esa ley y, \textbf{en caso de no aceptación por los propietarios, se deducen las correspondientes demandas de expropiación}.

El estándar de compensación no fue el valor de mercado: fue la valuación fiscal más treinta por ciento. Y el orden de los actos es el que define la experiencia de quien pierde la tierra: primero el Estado toma posesión depositando esa suma, después se discute el precio, y sólo si el propietario no acepta hay juicio.

Ese diseño explica por qué, medio siglo después, la memoria familiar de una expropiación puede diferir tanto del expediente. No hace falta que haya habido irregularidad: el propio mecanismo legal produce la sensación de despojo, porque quien pierde la tierra la pierde antes de haber discutido cuánto vale.

Y hay una asimetría de largo plazo que conviene enunciar sin adjetivos. La tierra se pagó a valuación fiscal más treinta por ciento en 1972. El agua que ese embalse almacena abastece hoy al sector norte de la ciudad de Salta. El municipio que la produce no tiene red de gas, tiene cobertura cloacal parcial y tiene desigualdad documentada en el acceso al agua potable, según el capítulo \ref{cap:redes}.

Una última precisión sobre el contexto institucional, porque cambia el sentido del acto.

La ley dice: ``El Gobernador de la provincia de Salta sanciona y promulga con fuerza de Ley''. No hubo Legislatura. En junio de 1972 la provincia estaba bajo gobierno de facto y las cámaras legislativas se encontraban disueltas; el Poder Ejecutivo sancionaba y promulgaba en un solo acto.

De modo que la decisión de expropiar seiscientas ochenta y seis hectáreas en el departamento La Caldera, con un estándar de compensación fijado por la propia norma y con toma de posesión previa a la discusión del precio, \textbf{no fue debatida por ningún cuerpo representativo}. Fue sancionada y promulgada el mismo día por la misma firma.

Es el antecedente más antiguo que este libro documenta de lo que sus dos tesis describen: una decisión sobre el suelo del departamento tomada sin deliberación pública, cuyos beneficios se realizan fuera del departamento y cuyos costos quedan adentro.

\begin{cautela}
Una advertencia sobre el aparato, que corresponde hacer antes de publicar.

La Ley 4486 invoca una «Ley provincial 4.172/68» y sus artículos 13 y 17. \textbf{Esa ley no existe con ese número.} La norma de 1968 que reformó esos artículos es la \textbf{Ley 4272}, sancionada y promulgada el 29 de noviembre de 1968 y publicada en el Boletín Oficial Nº 8.205 del 11 de diciembre, que modificó los artículos 13, 15 a 22, 35, 37 y 38 de la Ley 1336. Su nuevo artículo 17 es exactamente el estándar que se aplicó en Campo Alegre: el expropiante obtiene la posesión inmediata «consignando el importe de la última valuación fiscal \ldots\ incrementada en un 30\%». La cita de la Ley 4486 tiene, por lo tanto, un error de número. Lo que en cambio quedó establecido es que la ley orgánica sobre régimen de expropiación de la provincia de Salta es la \textbf{Ley 1336}, sancionada el 18 y promulgada el 24 de agosto de 1951, \textbf{renumerada como Ley 2614} por el Digesto; y que la Ley 2661, original 1383, de octubre de 1951, modificó su artículo 34. La correspondencia de numeración es la misma que este libro ya emplea para la tierra fiscal, donde la Ley 2616 es la original 1338.

Lo confirma la Ley 6354 de 1986 ---analizada más abajo---, que declara tramitar «por Leyes 1.336 y 4.272».

La afirmación de fondo no depende de la cita y está corroborada de manera independiente: el estándar de valuación fiscal incrementada en un treinta por ciento y la toma de posesión previa a la negociación del precio constan en el texto de las propias Leyes 4486 y 6181, que los aplican y los explicitan. Con el texto de la Ley 4272 a la vista, también la referencia al régimen general queda verificada.
\end{cautela}

\section{No fue la única}

La Ley 4486 no es un caso aislado. El digesto provincial registra al menos cuatro expropiaciones más en el departamento La Caldera, y su lectura conjunta matiza y precisa lo dicho más arriba.

\begin{ficha}{Ley provincial Nº 6181}
Promulgada el 20 de octubre de 1983. Boletín Oficial Nº 11.834 del 26 de octubre de 1983.
\begin{itemize}[leftmargin=*,nosep]
\item Declara de utilidad pública una hectárea del inmueble rural catastro Nº 132 del departamento La Caldera, plano de mensura y desmembramiento Nº 000248, cuyo dominio figura a nombre de \textbf{Hilario Arias}, con destino a la \textbf{Escuela Nº 660, Finca Mojotoro}.
\item Faculta a Fiscalía de Gobierno a iniciar el juicio de expropiación, liquidando la suma de \textbf{12,49 pesos argentinos, valor fiscal más el 30 por ciento}.
\item Dictada por el gobernador ``en ejercicio de las facultades legislativas concedidas por la Junta Militar''. Firmantes: Plaza --- Díaz --- Vicente --- Salazar --- Rodríguez.
\end{itemize}
\end{ficha}

\begin{ficha}{Ley provincial Nº 6354}
Expte.\ Nº 336-C/84. Sancionada el 4 de marzo de 1986, promulgada por Decreto Nº 884 del 21 de marzo. Boletín Oficial Nº 12.450 del 25 de abril de 1986.
\begin{itemize}[leftmargin=*,nosep]
\item Declara de utilidad pública el inmueble de propiedad de \textbf{Elena Serrey\index{Serrey, Elena} de González Bonorino y otros}, en la localidad de La Caldera, Sección A, Manzana 14, Parcela 1, \textbf{Matrícula Nº 1.303}, de 8.158,71 m².
\item Destino: \textbf{construcción de viviendas familiares}. Autoriza al Ejecutivo a donarlo al Instituto Provincial de Desarrollo Urbano y Vivienda una vez finiquitado el juicio.
\item Establece que el organismo beneficiario deberá iniciar la construcción en la totalidad del inmueble dentro de los dos años de la toma de posesión, \textbf{caso contrario volverá al dominio de la Provincia} sin interpelación y sin cargo.
\item Se tramita conforme la Ley de Expropiaciones Nº 1.336 y su modificatoria Nº 4.272.
\item Sancionada por la Legislatura en democracia. Firmantes: Canto --- Porrati --- Oliver --- Ulivarri.
\end{itemize}
\end{ficha}

\begin{ficha}{Leyes provinciales Nº 6334 y 6464 --- la Finca Vaqueros}
\textbf{Ley 6334} (sancionada el 3/9/1985; B.O.\ Nº 12.311): expropia \textbf{2 ha 1.736,80 m²} de la Finca Vaqueros, lote II~a, \textbf{catastro 1.229}, con dominio a nombre de \textbf{Manuel Serrey\index{Serrey, Manuel} y otros}. Destino: viviendas familiares, con donación al IPDUV y la misma cláusula de reversión a dos años que la Ley 6354. Se tramita por las Leyes 1336 y 4272.

\textbf{Ley 6464}: es un \textbf{decreto de necesidad y urgencia}, el 2876 del 14 de octubre de 1986, firmado en acuerdo general de ministros. Expropia \textbf{3.710,08 m²} del \textbf{catastro 1.280} de la misma finca, «de propiedad de la sucesión testamentaria de \textbf{Carlos Serrey}\index{Serrey, Carlos}», para la \textbf{planta potabilizadora} que debía dar agua corriente a Vaqueros. La Legislatura no lo aprobó ni lo rechazó: vencido el plazo del artículo 142 de la Constitución, el Decreto 1500 del 16 de julio de 1987 lo tuvo por ley (B.O.\ Nº 12.757).
\end{ficha}

\textbf{Son tres expropiaciones a titulares de apellido Serrey en dos años}, en los dos municipios del departamento: 1985 y 1986 en Vaqueros, 1986 en La Caldera. Dos fueron para vivienda social y una para agua potable. El apéndice \ref{cap:pedidos} pregunta si las viviendas se construyeron o si las tierras revirtieron a la Provincia.

\begin{cautela}
La primera consecuencia es una corrección de matiz sobre lo dicho en el apartado anterior.

El estándar de \textbf{valor fiscal más treinta por ciento} no fue un criterio adoptado especialmente para Campo Alegre en 1972: reaparece idéntico en 1983 para expropiar una hectárea destinada a una escuela rural. Era el régimen provincial ordinario de expropiaciones. Eso no lo vuelve generoso --- sigue siendo una compensación por debajo del valor de mercado y con posesión previa a la discusión del precio --- pero desactiva cualquier lectura que atribuya al caso Campo Alegre un trato singularmente desfavorable. Las familias de 1972 fueron tratadas como el régimen trataba a todos.

La Ley 6181 aporta además una cifra concreta que da la escala: \textbf{12,49 pesos argentinos por una hectárea}, valor fiscal más treinta por ciento, en 1983.

Hay una discrepancia normativa que conviene registrar. La Ley 4486 remite a la ``Ley Nº 4.172/68'' y la Ley 6354 a la ``Ley de Expropiaciones Nº 1.336 y su modificatoria Ley Nº 4.272''. No son dos regímenes: la primera cita tiene un error de número, porque la modificatoria de 1968 es la 4272 (véase la advertencia anterior).
\end{cautela}

\begin{cautela}
La segunda consecuencia es más importante y conecta con el capítulo \ref{cap:fincas}.

En 1986 el Estado provincial expropió tierra de \textbf{Elena Serrey\index{Serrey, Elena} de González Bonorino} en la localidad de La Caldera para construir viviendas familiares.

El capítulo \ref{cap:fincas} documentó que el mercado de suelo de Vaqueros nace de la resolución del juicio sucesorio de los descendientes de Serrey a principios de los años noventa, y que el edificio municipal de ese municipio se levantó en terreno donado por la sucesión del doctor Carlos Serrey\index{Serrey, Carlos}. El mismo apellido aparece también en el otro municipio del departamento, y en una posición distinta: como titular de un inmueble que el Estado debió expropiar para hacer vivienda social.

Los dos circuitos que este libro describe tienen, en este apellido, un punto de contacto documental. En Vaqueros, la partición hereditaria de un dominio Serrey\index{Serrey, Carlos} abrió el mercado privado de parcelas, y en 1985 y 1986 dos fracciones de la misma Finca Vaqueros pasaron al Estado por expropiación. En La Caldera, ocho mil metros cuadrados de un dominio Serrey pasaron al Estado por expropiación para vivienda familiar. Mismo apellido, mismo departamento, pocos años de distancia, mecanismos opuestos.

Las leyes de 1985 y 1986 nombran a Manuel Serrey y a la sucesión testamentaria de Carlos Serrey; el parentesco de Elena Serrey de González Bonorino con ellos no consta y requiere el antecedente sucesorio. No se afirma vínculo alguno entre los hechos más allá de esos nombres y la contigüidad territorial. Se consigna porque es exactamente la clase de coincidencia que el capítulo \ref{cap:fincas} planteó como hipótesis de larga duración, y que las leyes de expropiación vuelven documental.
\end{cautela}

La tercera observación es sobre el artículo 4º de la Ley 6354, que merece seguimiento propio.

La norma establece que si el Instituto Provincial de Desarrollo Urbano y Vivienda no inicia la construcción en la totalidad del inmueble dentro de los dos años de la toma de posesión, la tierra \textbf{vuelve automáticamente al dominio de la Provincia}. Es una cláusula de reversión por incumplimiento, poco habitual y estrictamente redactada.

Saber si esas viviendas se construyeron, y en su caso cuándo, o si la matrícula 1.303 revirtió al dominio provincial, es una pregunta con respuesta registral y con consecuencias sobre el stock de tierra fiscal disponible en el departamento --- que es el objeto del capítulo \ref{cap:tierrafiscal}.

\section{Campo Alegre antes del dique}

En 1972 un gobierno de facto declaró de utilidad pública y sujetas a expropiación las tierras
necesarias para construir el embalse Campo Alegre. Treinta y cuatro años antes, esa finca
aparece en el Boletín Oficial con dueño, título y folio.

\begin{ficha}{Edicto de concesión de agua --- B.O.\ Nº 1.746, junio de 1938}
Expediente 8312, letra F ---el decreto de 1939 lo cita como 8312-P/937; una de las dos letras es error de impresión o de lectura---. \textbf{Isaac Fernández}, comerciante, domiciliado en Caseros de la
ciudad de Salta, solicita concesión de agua del \textbf{río Santa Rufina} para riego ---
\textbf{veinticinco litros por segundo} sobre unas \textbf{doscientas hectáreas regables} ---
a favor de la finca \textbf{«Campo Alegre»}, departamento de La Caldera, \textbf{sin
mensurar}.

\textbf{Linderos:} «Norte, finca La Cañada Ancha, de Julio Mercado; Sud, finca El Angosto, de
\textbf{Daniel Linares\index{Linares, Daniel}}; Este, finca Loma Bola, de Julio Mercado; Oeste, inmueble Santa
Rufina, de \textbf{Bernardo Austerlitz\index{Austerlitz, Bernardo}}, hoy su heredera Carolina Austerlitz de Aráoz.»

\textbf{Título:} compra a Teodoro Sosaya, escritura del 12 de noviembre de 1937 ante el
escribano Carlos Figueroa, inscripta al \textbf{folio 374, asiento 475 del Libro D de Títulos
de La Caldera}.

La toma ocupa terreno de la sucesión Austerlitz\index{Austerlitz, Bernardo} y las acequias corren por Santa Rufina y por
Campo Alegre. Interviene la \textbf{Municipalidad de La Caldera}, cuyo informe funda el
decreto de trámite.
\end{ficha}

Cuatro cosas hacen de este edicto una pieza importante para el capítulo.

\textbf{Primera: el topónimo tiene dueño desde antes.} Campo Alegre es partido policial desde
1911 y finca con título inscripto desde 1937, comprada a Teodoro Sosaya por un comerciante de
la capital. Cuando la Ley 4486 lo declare de utilidad pública en 1972, no estará expropiando
un lugar sino un dominio con historia registral.

\textbf{Segunda: ya se pedía agua para esas tierras.} Veinticinco litros por segundo
del río Santa Rufina para regar doscientas hectáreas. La finca que treinta y cuatro años
después será el vaso del embalse ya tramitaba, en 1938, una
concesión de riego; el apartado siguiente muestra que le fue denegada.

\textbf{Tercera: los linderos son los de siempre.} Al sur, \textbf{Daniel Linares\index{Linares, Daniel}} --- el
presidente de la comisión de defensas del río de 1910, miembro del cuerpo municipal en 1922 y
propietario de Los Porongos en 1932 --- con la finca El Angosto. Al oeste, la sucesión de
\textbf{Bernardo Austerlitz\index{Austerlitz, Bernardo}}, propietario de San Alejo y Santa Rufina en 1925 y senador por el
departamento en 1927, ahora representada por su heredera. El apéndice \ref{cap:personas} reúne ambas
trayectorias.

\textbf{Y cuarta: la finca no estaba mensurada.} El edicto lo dice: «sin mensurar». Un
inmueble de al menos doscientas hectáreas regables, con título inscripto y concesión de agua
en trámite, cuyos límites nadie había medido. Es la misma condición que el capítulo \ref{cap:fincas}
documenta para las mensuras por mojones antiguos, nueve años después de la catastración
provincial de 1929.

\textbf{Y hay un eslabón anterior, que el barrido de 1914 incorporado al cierre puso a la
vista.} El edicto de 1938 deja la cadena de Campo Alegre empezando en 1937, con la compra de
Isaac Fernández a Teodoro Sosaya. Veintitrés años antes, la finca ya tenía dueño con el mismo
apellido y ya intentaba medirse.

\begin{ficha}{Deslinde de Campo Alegre --- B.O.\ Nº 512, 1914, h.\ 4}
«\textit{Habiéndose presentado el señor \textbf{Francisco Sosaya} con títulos bastantes
solicitando deslinde, mensura y amojonamiento de la finca \textbf{Campo Alegre} situada en el
departamento de la Caldera}», el juez de primera instancia en lo civil y comercial
\textbf{Francisco F.\ Sosa} ordena citar por treinta días a quienes se crean con derecho.
Agrimensor propuesto: \textbf{Ettore Chiostri}. Firma el secretario \textbf{Nolasco Zapata};
Salta, 16 de septiembre de 1914.

\textbf{Límites:} «\textit{Al norte, propiedad de los Mercado; al naciente, terrenos que
fueron de Abelino Costas; al sur, con los mismos sucesores de Costas y de José Fernández, y al
oeste, de estos mismos sucesores Fernández}».
\end{ficha}

Tres cosas se siguen de ahí, y la tercera es una advertencia.

\textbf{La primera: los Sosaya tenían Campo Alegre al menos desde 1914}, con «títulos
bastantes» que el edicto no describe. El comprador de 1937 le compra a un Teodoro Sosaya; el
solicitante de 1914 es un Francisco Sosaya, el mismo que en los avisos de San Alejo y Santa
Rufina figura como lindero del naciente (capítulo \ref{cap:siglo}). Qué relación hay entre los
dos, y si el título de 1937 desciende del deslinde de 1914, lo dice el antecedente dominial que
el apéndice \ref{cap:pedidos} ya pide, \textbf{y ahora puede pedirse con fecha y con juzgado}.

\textbf{La segunda: el deslinde de 1914 tampoco terminó en mensura}, o no la terminó de modo
que quedara asentada, porque el edicto de 1938 dice de la misma finca «sin mensurar».
Veinticuatro años entre el pedido y esa constancia.

\textbf{Y la tercera: los dos edictos no describen la misma figura.} En 1914 los linderos son
los Mercado al norte, los sucesores de Abelino Costas al naciente y al sur, y los sucesores
Fernández al sur y al oeste. En 1938 son La Cañada Ancha de Julio Mercado al norte, El Angosto
de Daniel Linares\index{Linares, Daniel} al sur, Loma Bola de Julio Mercado al este y Santa
Rufina de la sucesión Austerlitz\index{Austerlitz, Bernardo} al oeste. \textbf{Este libro no
concilia las dos}: pueden ser la misma finca con los vecinos cambiados en veinticuatro años, o
dos porciones distintas de un nombre que en 1911 era el de un partido policial entero. Sin
plano no se decide.

\subsection*{Y la concesión se denegó, por una ordenanza de 1914}

Un año después del edicto, el trámite terminó mal para el solicitante, y el decreto que lo
rechaza es el documento de mayor densidad institucional que este relevamiento encontró en
todo el período histórico.

\begin{ficha}{Decreto Nº 2805 del 24 de mayo de 1939 --- Expediente 8312-P/937}
«Visto el expediente \ldots\ en el cual Don Isaac Fernández solicita una concesión de agua del
Río Santa Rufina, para regadío del inmueble de su propiedad denominado «Campo Alegre»»; y
considerando:

«Que durante la tramitación del presente expediente se ha formulado \textbf{disenso municipal}
(fs.\ 19-20 y 69-72) y \textbf{oposiciones particulares} (fs.\ 26-28, 38-41 y 50-51),
aduciéndose como principal fundamento la existencia de una \textbf{Ordenanza dictada por el
Municipio de La Caldera en fecha octubre 13 de 1914, que reglamenta la distribución de la
totalidad del caudal de agua del Río Santa Rufina}»;

«Que no obstante constituir una facultad propia del Poder Ejecutivo el otorgamiento de
concesiones de uso de agua para regadío (artículo 112 del Código Rural), en el presente caso
existe una Ordenanza aprobada por el \ldots\ que \textbf{reconoce los derechos adquiridos por
uso tradicional}»;

«Que de acordarse la concesión solicitada produciría como consecuencia inmediata \textbf{la
alteración de la forma en que se distribuye el agua} del referido río».

\textbf{Art.\ 1º: Desestímase el pedido.} Firman Luis Patrón Costas y Carlos Gómez Rincón.
\end{ficha}

Cuatro cosas.

\textbf{Primera: existía una ordenanza municipal de 1914 que repartía todo el caudal de un
río.} No una parte: «la totalidad». El municipio de La Caldera había reglamentado, seis años
después de aparecer el Boletín Oficial y con una renta anual de poco más de mil pesos, la
distribución íntegra del agua del Santa Rufina.

Este libro no conoce esa ordenanza --- no aparece en \textbf{ninguna de las sesenta y nueve
ediciones de 1914} que este relevamiento leyó una por una, con reconocimiento de texto propio
sobre las 279 hojas del año, ni en ningún repositorio consultado; las siete ediciones que el
repositorio no tiene son de junio y la ordenanza es del 13 de octubre (apéndice
\ref{cap:fuentes}) ---
y su ausencia es notable: \textbf{es el reglamento de reparto del agua más antiguo del departamento
que el archivo menciona, y este relevamiento no encontró su texto.} (Más antigua es la
concesión de 1885 del propio municipio a Eustaquio Murúa, pero es un acto singular, no un
reparto.)

\textbf{Y el propio decreto dice dónde buscarla.} El considerando sitúa el disenso municipal
en las \textbf{fojas 19-20 y 69-72} del expediente. Para oponerse invocando una ordenanza de
veinticuatro años antes, el municipio tuvo que acompañarla o transcribirla: \textbf{la copia
más probable de la ordenanza de 1914 está en esas fojas}, en un expediente provincial, y no en
La Caldera. Es el mismo movimiento que este libro ya hizo una vez en sentido inverso, cuando el
reglamento que el Boletín no publicó apareció en el digesto del municipio que lo había dictado
(capítulo \ref{cap:vaqueros}): \textbf{el documento no está donde lo busca quien lo busca,
sino donde lo necesitó quien lo usó}.

\begin{cautela}
Conviene no trasladar sin más el caso de Vaqueros a éste. Vaqueros tiene un digesto en línea
desde 2024 y La Caldera no publica sus ordenanzas de manera accesible (capítulo
\ref{cap:metodo}). Y en 1914 el municipio no tenía Concejo ni boletín propio: lo gobernaba una
comisión municipal designada por la Provincia, con una renta anual de poco más de mil pesos
(capítulo \ref{cap:siglo}). Lo que puede existir de esa época en el municipio es un
\textbf{libro de actas}, no una colección de ordenanzas publicadas. Este libro no sabe si se
conserva.
\end{cautela}

\textbf{Segunda: esa ordenanza reconocía «derechos adquiridos por uso tradicional».} Es decir:
en 1914 el municipio no inventó un reparto sino que registró uno anterior, consuetudinario.
Había, antes de toda concesión provincial, un orden de usos del agua que la comunidad se había
dado y que la norma se limitó a poner por escrito.

\textbf{Tercera: la Provincia se abstuvo de pasar por encima.} El decreto reconoce que otorgar
concesiones es facultad propia del Poder Ejecutivo, y aun así deniega, porque hacerlo
alteraría el reparto municipal. \textbf{En 1939, una ordenanza local sobre agua prevaleció
sobre una solicitud de concesión provincial.}

Conviene medir esa afirmación contra el resto del libro. El capítulo \ref{cap:ribera} documenta veinticinco
actos provinciales de línea de ribera entre 2013 y 2015 en los que no consta intervención municipal alguna ---la excepción, contada aparte, es la Resolución 023/14, pedida por el intendente sobre el propio pueblo---; el capítulo \ref{cap:loteo} documenta que el plano de un loteo se aprobó por resolución
ministerial sin que figure ninguno de los diez actos que el Código municipal exige. En
1939 el municipio tenía voz en el agua de su territorio, y la ejerció.

\textbf{Y cuarta: hubo oposiciones de particulares.} Tres presentaciones distintas, en fojas
identificadas. Vecinos o regantes se opusieron a que un comerciante de la capital tomara agua
del río, y ganaron. Con el reclamo de Campo Santo de 1928 contra la concesión de Urquiza (capítulo \ref{cap:politica}), es uno de los dos casos del período, entre los que registran las ediciones leídas, en que una oposición a un trámite de agua prospera, y el único que se sostuvo: aquella revocación la anuló la Corte en 1931.

Fernández pidió reconsideración, y en enero de 1940 el Poder Ejecutivo confirmó la
denegatoria con un fundamento más fuerte que el anterior: \textbf{razones de orden legal
impiden al Poder Ejecutivo otorgar concesiones sobre el agua del río Santa Rufina}. No que
no correspondiera en ese caso: que no podía.

De modo que en 1940 el río Santa Rufina estaba, en los hechos, \textbf{fuera del alcance del
poder concesional de la Provincia}, repartido íntegramente por una ordenanza municipal de
1914 que reconocía usos anteriores a ella. Es la situación exactamente inversa a la que este
libro documenta para el presente.

\begin{cautela}
Esto obliga a corregir una impresión que el apartado anterior podía dejar. \textbf{No consta que la finca
Campo Alegre tuviera derecho de agua reconocido cuando se la expropió en 1972}: lo había
solicitado en 1937 y se lo denegaron en 1939, y el período posterior a 1945 este libro no lo leyó de corrido.

Lo que existía sobre el río Santa Rufina era otra cosa, y más interesante: un reparto
consuetudinario entre usuarios, puesto por escrito en una ordenanza municipal de 1914 y
sostenido por la Provincia contra un pedido de concesión nuevo. Qué pasó con esos derechos
de uso tradicional cuando se construyó el embalse es una pregunta que este libro no puede
contestar y que ninguna de sus fuentes menciona.
\end{cautela}

\begin{cautela}
El edicto entrega algo que el resto del archivo casi nunca da: \textbf{el folio, el asiento y
el libro de la inscripción}, más el nombre del vendedor, la fecha de la escritura y el
escribano. Es, junto con el folio 342 del Libro B correspondiente a la finca Chalchanio, la
segunda referencia registral concreta que este relevamiento obtuvo.

Con esos dos datos --- folio 374, asiento 475, Libro D --- se puede empezar a reconstruir la
cadena de titularidad de Campo Alegre hasta la expropiación de 1972. Es un pedido de una
línea en el apéndice \ref{cap:pedidos}, y contestaría quién era dueño de la tierra del embalse cuando el
Estado decidió tomarla.
\end{cautela}

\section{Ocho años después, el dique tiene nombre}

La Ley 4486 expropió en 1972 las tierras para el embalse. En 1980 la obra ya está en pie,
tiene nombre propio y sigue recibiendo presupuesto.

El Plan Analítico de Trabajos Públicos de ese año, aprobado por la Resolución 163 del
Ministerio de Economía, incluye dos obras de la Administración General de Aguas de Salta en
el \textbf{Dique Ing.\ José Alfonso Peralta}, La Caldera: el \textbf{proyecto del
establecimiento potabilizador}, por \textbf{trescientos millones de pesos} con cargo a Rentas
Generales, y la \textbf{toma y canal de aducción} del embalse.

Dos observaciones.

\textbf{El embalse lleva el nombre del ingeniero que presidía Vialidad en 1933.} El capítulo \ref{cap:hacienda}
documenta el acta del Directorio que ese año le encomendó al presidente municipal de La
Caldera la limpieza de un camino: la presidía el ingeniero \textbf{José Alfonso Peralta}.
Este libro no puede establecer si es la misma persona --- haría falta el acto que impuso el
nombre al dique ---, pero el nombre completo coincide y el cargo es del mismo ramo.

\textbf{Y la potabilizadora es de 1980.} Trescientos millones de pesos para proyectar el
establecimiento que trata el agua del embalse. Cuarenta y seis años después, el presidente
del centro vecinal de La Calderilla declara que esa planta está en el departamento, que hay
cavernas subterráneas que captan agua por filtración debajo del lecho del río, y que los
loteos están aguas arriba de esas tomas.

\section{Lo que el acueducto se llevó}

Este capítulo documenta que en 1972 se expropiaron seiscientas ochenta y seis hectáreas para
el embalse y que el acueducto que sale de él abastece a la ciudad de Salta. El diagnóstico
social del Plan de Manejo agrega qué costó esa obra en el terreno, y lo dicen los vecinos.

\begin{ficha}{Plan de Manejo, página 134 --- desmonte y bosque ribereño}
El relevamiento consigna el reclamo «desde la \textbf{eliminación de un bosque ribereño
durante la construcción del acueducto Campo Alegre}». Se avistaron carteles de un \textbf{plan
de forestación compensatoria}, y por esa razón «en un taller anterior a este trabajo la
población había solicitado inventarios forestales en la costanera del río». A campo se
constataron arbustos --- guarán, palo bobo, churquis y tuscas --- y «algunos ejemplares de
sauce».

Testimonios recogidos, entre comillas en el original:

«Durante los últimos años, \textbf{se ha loteado mucho, no han respetado nada} y las
pendientes hacen que el agua corra como río en el verano.»

«Cuando se hizo el acueducto \textbf{se desmontó todo el bosque del borde del río, eso cuidaba
para que el agua no pasara}, pero ahora se lleva puesto todo lo que encuentra. ¿Qué hacen los
municipales? Cortan los pocos árboles que quedan para ponerlos en la orilla, más daño que
beneficio, \textbf{culpa de eso el río casi se lleva el cementerio}. Ni los muertitos se van a
salvar.»

«Desmontes y falta de legislación que controle \textbf{sobre todo a los nuevos loteos}.»
\end{ficha}

\textbf{Primera: la obra que lleva agua a la capital se llevó la defensa natural del pueblo.}
El bosque ribereño es, en términos hidráulicos, lo que frena la erosión de márgenes. El capítulo \ref{cap:defensas} documenta ciento dieciséis años de defensas de rama, piedra y hormigón construidas
contra ese río. \textbf{Una parte del problema que esas obras combaten la produjo otra obra
del mismo Estado.}

\textbf{Segunda: hubo forestación compensatoria, y quedan carteles.} No se sabe si se
ejecutó ni en qué medida; lo que el relevamiento constató a campo son arbustos y «algunos
ejemplares de sauce». Y la población había pedido, en un taller anterior, inventarios
forestales de la costanera.

\textbf{Y tercera, que es la que más dice sobre el capítulo \ref{cap:plan} de este libro.} El Plan de
Manejo asigna doscientos veintiséis millones de pesos a «bosques de ribera» --- menos del uno
por ciento de su inversión total. \textbf{Es dinero para reponer lo que otra obra pública
eliminó}, y es la única medida de infraestructura verde del cuadro.

\begin{cautela}
Las tres citas son de vecinos anónimos recogidas por la consultora, no de un informe técnico.
Este libro no puede establecer cuánto bosque se eliminó, en qué año exacto, con qué
autorización ni qué alcance tuvo la forestación compensatoria.

Lo que registra es que \textbf{el reclamo existe, está documentado en un informe oficial y
atribuye a la obra del acueducto la pérdida de la defensa natural del río}. Y que la
existencia de carteles de forestación compensatoria confirma, al menos, que hubo una
obligación de reponer.
\end{cautela}

\subsection*{La traza del acueducto, con coordenadas}

El expediente del amparo incorporó, como antecedente, un estudio de ruido ambiental
realizado por \textbf{Aguas del Norte} los días 7 y 8 de julio de 2015 sobre la «Planta La
Caldera y traza del Acueducto Norte». De sus diez puntos de medición, la Secretaría de
Ambiente retomó cinco, con sus coordenadas:

\needspace{6\baselineskip}
{\small
\begin{longtable}{@{}ll@{}}
\toprule
\textbf{Punto} & \textbf{Coordenada} \\
\midrule
\endhead
R01 & 24°35'39,10"S --- 65°22'32,40"O \\
R02 & 24°36'9,10"S --- 65°22'35,20"O \\
R03 & 24°36'37,70"S --- 65°22'54,30"O \\
R04 & 24°36'59,50"S --- 65°23'15,20"O \\
R05 & 24°41'30,75"S --- 65°23'41,85"O \\
\bottomrule
\end{longtable}
}

Los cuatro primeros se escalonan a lo largo de unos dos kilómetros y medio en sentido
norte-sur; el quinto está unos ocho kilómetros y medio más al sur, ya en dirección a Vaqueros.
\textbf{Es la primera vez que este libro puede ubicar la infraestructura del acueducto sobre
el terreno con precisión de segundos.}

\begin{cautela}
\textbf{Y esos cinco puntos admiten un cruce que este libro deja planteado y no resuelve.} El
capítulo \ref{cap:ribera} transcribe las coordenadas Gauss-Krüger de la línea de ribera que la
Resolución 185/14 fijó sobre el río La Caldera en la matrícula 4366 ---la finca que el edicto de agua de 2015 vincula con
el loteo El Durazno, aprobado sobre la matrícula 6.558---, y el capítulo \ref{cap:tierrafiscal} muestra que los dos sistemas son
convertibles: con los puntos del plano 000962 se estableció así la distancia al paraje El
Gallinato.

Hecha esa conversión de manera aproximada, los dos extremos del tramo de ribera de la 185/14
caen \textbf{dentro del segmento que forman los puntos R04 y R05} de este estudio, a unos pocos
cientos de metros de la recta que los une. \textbf{Este libro no lo afirma}: la conversión
esférica que permite hacer una planilla tiene un error del mismo orden que la distancia medida,
los cinco puntos son estaciones de medición de ruido y no el eje de la cañería, y la convención
de las columnas «X» e «Y» del edicto se transcribe tal como aparece y no está verificada.

Lo que sí queda establecido es que \textbf{la comprobación es posible y barata}, porque las dos
series ya están transcriptas en este libro: reproyectarlas correctamente contesta si la traza
del acueducto corre o no sobre la franja ribereña que los vecinos denunciaron en 2016, que es
la pregunta que el apartado siguiente deja abierta. \pendiente{reproyección en coordenadas
planas de los puntos del estudio de ruido y de las coordenadas de la Resolución 185/14, y traza
oficial del Acueducto Norte}
\end{cautela}

Y conviene retener de dónde sale el dato: no de un acto de la empresa ni del organismo
hídrico, sino \textbf{de un estudio de ruido incorporado a un expediente judicial de amparo}.
El capítulo \ref{cap:opacidad} describe esa forma de acceso: la información existe, no está negada, y aparece
por una vía lateral que nadie diseñó para eso.

\section{La emergencia hídrica de 2013}

Cinco años antes de los aluviones que originan el amparo colectivo, la Provincia había declarado formalmente que se estaba quedando sin agua.

\begin{ficha}{Ley provincial Nº 7781 --- Decreto de Necesidad y Urgencia Nº 1595/13}
Dictado el 30 de mayo de 2013, promulgado el 30 de agosto, publicado en el Boletín Oficial Nº 19.138 del 3 de septiembre de 2013. Convertido en ley por el transcurso del plazo del artículo 145 de la Constitución Provincial sin pronunciamiento de las Cámaras.
\begin{itemize}[leftmargin=*,nosep]
\item Declara el \textbf{estado de Emergencia Hídrica por escasez de agua en todo el territorio de la Provincia}, por el término de un año.
\item Fundamento, según sus considerandos: la crisis obedece fundamentalmente al \textbf{cambio climático}, que provoca merma sustancial de agua en las cuencas, diques, ríos y arroyos, y al déficit de precipitaciones de los últimos seis meses.
\item Consigna que la falta del recurso incide en el suministro de agua potable a comunidades, establecimientos educativos y grandes centros urbanos, y en la producción agrícola-ganadera.
\item Ordena racionalizar el uso del agua y \textbf{readecuar y ordenar todas las cuencas}.
\item Origen: Resolución Nº 142/13 de la Secretaría de Recursos Hídricos.
\end{itemize}
\end{ficha}

Este decreto importa por dos razones.

\textbf{Primera.} En 2013 el Poder Ejecutivo provincial ya atribuía la crisis hídrica al cambio climático, en un instrumento con fuerza de ley. Ocho años después, la sentencia del amparo colectivo introdujo el cambio climático como hecho notorio y citó la Opinión Consultiva OC-23/2017 de la Corte Interamericana. No fue el tribunal quien trajo el argumento climático a la discusión provincial: el Ejecutivo lo había puesto en un decreto de necesidad y urgencia.

\textbf{Segunda, y es la que ordena el capítulo.} El decreto ordena \textit{readecuar y ordenar todas las cuencas}, y la Secretaría de Recursos Hídricos es el organismo que lo propuso. Seis años después, esa misma Secretaría, citada en el amparo por los aluviones de la microcuenca del río La Caldera, pidió el rechazo de la demanda cuestionando la legitimación del actor y sosteniendo que el organismo no tenía capacidad para hacer obras.

El mismo organismo que en 2013 aconsejó al Ejecutivo declarar la emergencia y ordenar las cuencas, en 2019 declaró ante un tribunal que no podía hacerlo.

\section{Lo que se construyó sobre esa tierra}

El embalse Campo Alegre --- denominación oficial \textit{Ing.\ José Alfonso Peralta} --- es un lago artificial situado a unos cinco kilómetros del pueblo de La Caldera y a unos veinticinco de la capital, sobre el curso del río La Caldera y a la vera de la Ruta Nacional 9. Se comenzó a construir en 1972, el mismo año de la ley; la adjudicataria fue Rodio Argentina S.A. Tiene cuarenta y ocho hectómetros cúbicos de capacidad, trescientas veintidós hectáreas de espejo y cuarenta metros de profundidad máxima, a 1.421 metros sobre el nivel del mar. Tardó más de dos años en llenarse.

Fue concebido como reserva de agua potable para el sector norte de Salta capital --- incluyendo a Vaqueros y La Caldera --- y como fuente de riego.

La planta potabilizadora ubicada en La Caldera recibe el agua en una cisterna principal, desde la cual un acueducto de aproximadamente veinte kilómetros ---veintidós, según la inauguración oficial--- atraviesa ambos municipios hasta el nudo de distribución de El Huaico, en avenida Bolivia y Patrón Costas. La primera etapa se concluyó hacia 2021, con financiamiento de la Corporación Andina de Fomento.

Conviene poner las dos fechas juntas. La tierra se tomó en 1972, con posesión inmediata y compensación fijada en valuación fiscal más treinta por ciento. La obra que la vuelve útil para la capital --- el acueducto de unos veinte kilómetros hasta El Huaico --- se completó en su primera etapa en 2021, con financiamiento internacional.

Cuarenta y nueve años entre una cosa y la otra. En el medio, el departamento pasó de tener menos de mil habitantes en su cabecera a ser el cuarto de mayor crecimiento demográfico del país, perdió el 88 por ciento de su superficie agropecuaria declarada ---sobre todo monte y pastizal--- y su municipio cabecera no recibió red de gas ni cobertura cloacal completa.

\begin{cautela}
\textbf{Una precisión sobre qué agua va por dónde, porque las fuentes no dicen lo mismo.} La planta del dique y su acueducto de veintidós kilómetros se inauguraron oficialmente en octubre de 2021 como «Planta Potabilizadora Dique Campo Alegre y Acueducto Norte», con una producción de cuatro mil metros cúbicos por hora tomada del embalse; la segunda etapa se licitó por cincuenta millones de dólares. El informe que Aguas del Norte elevó a Fiscalía de Estado en diciembre de 2022 (capítulo \ref{cap:amparo}) describe, en cambio, un \textbf{Acueducto Norte} que capta del río a la altura de la curva de El Gallinato ---la captación cuyas cámaras se protegieron en 2014 (capítulo \ref{cap:defensas})---, un \textbf{Acueducto Wierna} que capta en el río Las Nieves y un \textbf{Acueducto Campo Alegre} al que a esa fecha le faltaban «detalles menores» para operar. Las dos versiones pueden conciliarse ---el nombre designa sistemas de épocas distintas, y una obra inaugurada puede no estar todavía en operación regular---, pero este libro no puede establecer cuál es la situación actual. Qué parte del agua que recibe la capital sale hoy del embalse y qué parte del río es un dato que el apéndice \ref{cap:pedidos} pide.
\end{cautela}

El agua salió. Los servicios no entraron.

\subsection*{Y en 2019, una turbina al pie de la presa}

\begin{ficha}{Concesión de agua pública --- B.O.\ Nº 20.451 a 20.455, 25 de febrero a 1 de marzo de 2019}
Expediente 0090034-179036/2018. La sociedad \textbf{Skru S.A.}, por su director, gestiona la concesión para \textbf{aprovechamiento de energía hidráulica} con una potencia instalada de \textbf{907 kW}, «para el proyecto a localizarse en el \textbf{pie de la presa dique Campo Alegre} mediante una turbina de flujo», en el catastro 1782 del departamento. Concesión «real y permanente» (art.\ 92 del Código de Aguas). Oposiciones en treinta días hábiles. Edicto del 21 de febrero de 2019.
\end{ficha}

\begin{cautela}
El embalse construido sobre la tierra expropiada en 1972 para abastecer de agua a la capital aparece, casi medio siglo después, como emplazamiento de un proyecto privado de generación. El Código de Aguas lo prevé: la energía hidráulica es un uso concesionable, y una turbina al pie de una presa aprovecha un caudal que ya se erogaba. El edicto no dice si el proyecto se concesionó ni si se construyó, ni qué relación tiene el catastro 1782 con el inmueble expropiado. Lo que sí deja escrito es que el mismo caudal que sale del dique hacia la planta potabilizadora tuvo, en 2019, un tercer interesado además del Estado y de los regantes. \pendiente{resolución recaída en el expediente 0090034-179036/2018, contrato o permiso de uso de la presa, y estado del proyecto de Skru S.A.}
\end{cautela}

\section{La advertencia de 2016}

En 2016, un grupo de vecinos autoconvocados de La Caldera tomó conocimiento del proyecto del Sistema Acueducto Norte y advirtió que su ejecución implicaría el desmonte del cordón ribereño de árboles entre la costanera y el río. La movilización buscaba evitar tanto la afectación del bosque ribereño como la exposición de la población a inundaciones.

En febrero de 2019, en el marco de la crisis ambiental del departamento, los denunciantes cuestionaron ante la prensa la ausencia de monitoreo efectivo sobre la traza y construcción del acueducto, y sostuvieron que no era aconsejable su instalación sobre una franja ribereña de alta fragilidad hidrológica cercana al único acceso a La Caldera, ni sobre la zona donde se encontraba un bosque ribereño nativo de alta biodiversidad que fue desmontado parcialmente para la traza.

\begin{cautela}
El orden de los hechos importa más que su contenido individual. Vecinos del mismo pueblo advirtieron en 2016 sobre el desmonte del bosque ribereño para el acueducto; en diciembre de 2018 y enero de 2019 ocurrieron los aluviones que originan el amparo colectivo; y en 2019 los denunciantes volvieron sobre la falta de control de esa misma traza. Este libro no establece que los autoconvocados de 2016 y los damnificados de 2018 sean las mismas personas.

La secuencia no prueba causalidad --- eso requiere peritaje hidrológico, no cronología --- pero establece algo relevante: la advertencia fue formulada públicamente, por vecinos organizados, dos años antes del daño. Cualquier defensa institucional basada en la imprevisibilidad del evento tiene ese antecedente en contra, y se suma a la cota máxima edificable que el propio plan de 2015 dejó como pendiente.
\end{cautela}

\section{La contradicción de fondo}

El departamento produce el agua de la capital y remueve, para lotear, el bosque de la cuenca que la produce. Simultáneamente, en Vaqueros, los vecinos históricos de la zona alta reclaman desde hace más de quince años acceso al agua potable en sus domicilios, mientras los emprendimientos privados que se instalan en la misma zona resuelven su abastecimiento por cuenta propia.

En julio de 2021 se manifestaron reclamando acceso igualitario. Según \textit{Página/12}, denunciaron que se priorizaba el acceso rápido de los emprendimientos inmobiliarios mientras ellos llevaban más de quince años de gestiones. Una vecina oriunda del pueblo relató que en la zona alta no tienen agua en las casas, que deben comprarla o esperar el camión municipal dos o tres veces por semana. Otro vecino, residente desde 2007, señaló que el lugar se fue poblando sin agua ni cloacas para los primeros habitantes, y que ya en 2010 comenzaron nuevos proyectos de barrios privados --- mencionó el loteo Las Vertientes, que cuenta con suministro desde las vertientes cercanas que le dan nombre. Observó que Vaqueros está rodeado de cuencas --- los ríos Yacones, De las Nieves, Vaqueros y Wierna --- y que por lo tanto el recurso existe: lo que hay es desigualdad en su distribución. Los manifestantes aclararon que no se oponen a la construcción de nuevas viviendas, sino a que no haya solución para los radicados y nacidos en el pueblo.

Aquí H.20 deja de ser hipótesis abstracta y se vuelve descripción de un caudal. No hay escasez hídrica en Vaqueros: hay un orden de prelación. Un emprendimiento nuevo resuelve su abastecimiento desde una vertiente propia en el plazo de su desarrollo comercial; un vecino histórico espera el camión municipal dos o tres veces por semana después de quince años.

La formulación estructural correcta no es que alguien decidió postergarlos. Es que el suelo con servicios se produce por inversión privada, y la inversión privada produce servicios para su propia parcela. El Estado municipal no arbitra ese orden porque no tiene con qué, y el capítulo \ref{cap:hacienda} documenta con qué exactamente no cuenta.

\section{El perilago}

En enero de 2026, se promulgó la Ordenanza Nº 830/25, sancionada por el Concejo Deliberante, que regula el ingreso al perilago del embalse y prohíbe expresamente todo cobro para acceder al espejo de agua. La iniciativa fue impulsada por el presidente del cuerpo, Germán Sumbay\index{Sumbay, Germán}, quien según declaraciones a un programa radial respondió a denuncias por el cobro de un peaje informal de entre dos mil y tres mil pesos para circular por caminos vecinales y acceder al dique, principalmente en el sector norte, sin contraprestación de servicios. La ordenanza establece que los campings privados sólo pueden cobrar por el uso efectivo de sus servicios, no por el acceso al embalse, por tratarse de administración estatal; dispone señalización oficial, intervención de las fuerzas de seguridad y apertura de un nuevo acceso público.

El caso del perilago es la contracara del caso Solanas que documenta el capítulo \ref{cap:loteo}, y por eso conviene retenerlo. En una misma semana, el mismo cuerpo deliberativo pide la clausura del ingreso a un loteo cerrado y sanciona una norma para abrir el acceso público a un bien del Estado que estaba siendo cobrado por privados.

La conclusión es incómoda para cualquier tesis simple: en el departamento La Caldera la disputa por el suelo y el agua no enfrenta a un poder local capturado contra una sociedad civil indefensa. Atraviesa al poder local por dentro.
